\section{Evidence Contract}
\label{app:evidence-contract}

An empirical claim may replace an \texttt{EVIDENCE BLOCKED} marker only when its
artifact chain binds: (i) the immutable preregistration and amendment history;
(ii) model, inference, runtime, grader, task-group, and price-table hashes;
(iii) every configured search closure, including failed calls and retries;
(iv) the signed registry of committed edits and one-time sealed aggregate
release; (v) the
prespecified estimate, uncertainty interval, and multiplicity result; and (vi)
the deterministic table or figure source.  A result exposed before these
conditions is development evidence and stays out of the confirmatory table.

\section{Primary Hypotheses}
\label{app:hypotheses}

The commit-reliability family contains two one-sided hypotheses:
\begin{align*}
H_{0R1} &: p_{\mathrm{FP}}(\text{\fullarm}) \geq 0.05,
  && H_{1R1}: p_{\mathrm{FP}}(\text{\fullarm}) < 0.05,\\
H_{0R2} &: p_{\mathrm{FP}}(\text{\scorearm})
 -p_{\mathrm{FP}}(\text{\fullarm}) \leq \delta_{\mathrm{FP}},
  && H_{1R2}: p_{\mathrm{FP}}(\text{\scorearm})
 -p_{\mathrm{FP}}(\text{\fullarm}) > \delta_{\mathrm{FP}}.
\end{align*}
Here \(\delta_{\mathrm{FP}}>0\) is justified from a signed operational risk
budget and frozen before the outcome-bearing pilot.  The
sealed-usefulness family tests
\(H_{0U1}:\Delta_{\mathrm{policy}}\leq0.02\) and
\(H_{0U2}:\Delta_{\mathrm{e2e}}\leq0.02\) against larger effects.  Each family
uses one-sided Holm--Bonferroni FWER 0.05, and every null in both families must
be rejected.  The full headline additionally requires simultaneous protected-
behavior and cost non-inferiority plus provenance, leakage, and missingness
closure.  The conjunction prevents a reject-all policy from passing on low
false promotion alone.  Verified-repair improvement at the prospective 0.15
margin, harmful promotion, and accepted-edit precision are key secondary
families rather than substitutes for this decision.
The 0.02 utility, 0.15 repair, protected-behavior, and cost margins likewise
require operational justification and are fixed before that pilot; pilot effect
sizes are nuisance/power inputs, not margin selectors.

Any ten-percentage-point statement is an independent preregistered headline
family on the binary-success scale.  Its exact contrasts and simultaneous
multiplicity adjustment must be fixed and powered before search; every one-
sided lower bound must be strictly greater than 0.10.  A normalized score does
not acquire percentage-point units.  The family cannot be declared after
unsealing, and sealed outcomes cannot trigger sample-size continuation or
method iteration in the same lineage.

\section{Prospective Sizing Boundary}
\label{app:power-proxy}

The repository currently has two method-development paths.  The bounded
analytic sensitivity proxy uses a latent Gaussian rank-copula parameter,
independent cell aggregates, a continuous approximation to verified repair,
and one non-conservative plug-in failure scenario.  A newer fixed-roster
prototype computes the equal-cell macro estimand and a basic nested bootstrap
over search seeds, task groups, and repeats, and passes synthetic datasets
through the same prototype analyzer during prospective simulation.  Neither
path has completed coverage calibration, frozen an authority-attested input
contract, or been adopted as the confirmatory estimator.  Neither can select a
formal seed count or produce confirmatory evidence.

Calibration-null, independent audit-null, and alternative streams are kept
separate where implemented, but their FWER-like, rejection, and Monte Carlo
rates remain diagnostics rather than coverage claims.  Caller-declared digests
or dataset-source labels do not verify artifact existence, observation status,
pilot disjointness, sealed closure, signatures, or blinding.  The current
selector's \emph{combined sizing event} is the conjunction of all three
Holm-corrected SESOI-margin rejections and the two real-task simultaneous
one-sided lower bounds above 0.10.  It excludes PURE, PURE@B, and external
claim-family sizing as well as protected-behavior, cost, provenance, and leakage
constraints.  More importantly, it contains no family-level false-promotion
endpoint or harmful-promotion process and therefore predates the present
co-primary design.  It is not the full headline pass and cannot freeze the
final sample size.  At a true effect exactly on its minimum boundary, a valid
one-sided margin test rejects only at its calibrated type-I rate; 0.90 combined-event
power requires a prospectively justified alternative strictly above every
margin.

The prototype generates its legacy repair endpoint independently of its two
real-task endpoints and simulates no false-promotion, failure, or missingness
process.  Its diagnostic null stream covers only the legacy joint SESOI
boundary---not partial nulls, reliability nulls, or 0.10 boundaries---and
neither proves that boundary least favourable nor gates selection.  It chooses the first
candidate whose 95\% Wilson lower bound for calibration combined-event power
meets a caller-declared target, then applies the criterion once on an independent
validation stream; failure yields no recommendation and no larger fallback.  A
formal design must freeze the target and roster; simulate family-level false and
harmful promotion, plausible endpoint dependence, ICC, and failure modes; and
make independent boundary/partial-null strong-FWER audits a fail-closed
prerequisite.

\section{Detailed Evaluation Protocol}
\label{app:evaluation-details}

\paragraph{Feedback and budget closure.}
\scorearm{} and \fullarm{} receive the same authenticated aggregate score,
interval, cost, and performance-only decision over utility, protected behavior,
and cost.  SCORE cannot read item outputs, trajectories, localization, or any
activation, adherence, behavior, revert, or placebo packet.  FULL may read only
the source-linked mechanism fields authorized before search.  A confidence
projection derives its level, endpoints, estimator version, and analysis-plan
identity from one authority-owned objective artifact; callers cannot supply
interval endpoints.  Every solver, diagnoser, proposer, materializer, ranker,
nominator, mechanism-judge, grader, retry, and failed-call charge remains in
the ledger.  Matching fixes ex-ante coordinate
rules and ceilings, not post-treatment candidate contents or realized tokens.

\paragraph{References and baselines.}
\purearm{} receives the official task payload plus provider-required
serialization, with no system prompt, examples, memory, skill, tool,
middleware, or routing logic.  \(\textsc{Pure@B}\) reserves FULL's full
model-call and token ceiling before execution and maps its complete independent
sample multiset to one output through a total, task-native, nonadaptive rule.
Random-valid uses a frozen finite catalog; prompt-only cannot serve as SCORE if
it receives item feedback.  Meta-Harness-style and Self-Harness optimizers are
secondary: missing public components make an implementation a paper-algorithm
reconstruction, and any stronger role model makes it heterogeneous and
ineligible for same-model or ten-point claims.

\paragraph{Tasks, sealing, and failures.}
Availability-only smoke tests select endpoint feasibility without benchmark
scores.  Tasks bind source/template/generator families, binary rubrics, group
weights, and a never-opened sealed partition before search.  All conditions and
failed runs reach terminal closure before one authorization evaluates every
committed-edit hash and final champion.  Malformed output, refusal, candidate-
caused timeout, or crash scores zero with cost retained.  A preregistered
infrastructure retry replaces the complete matched block, not one favorable
arm; revision drift invalidates the entire model--task--seed block.  The primary
analysis is intention-to-treat with failure-as-zero, complete-case and worst-
case sensitivity analyses.  We report absolute outcomes, protected slices,
uncached/cached input, output/reasoning tokens, normalized and billed cost, wall
time, failures, retries, score per million tokens, and a Pareto frontier.

\section{Prospective HarnessFaultBench Design}
\label{app:fault-design}

\begin{table}[h]
  \caption{Prospective fault factors.  Semantically impossible combinations are
  compatibility-masked before sampling.}
  \label{tab:fault-factors}
  \centering
  \begin{tabular}{p{1.25in}p{3.80in}}
    \toprule
    \textbf{Factor} & \textbf{Levels} \\
    \midrule
    Mutable surface & prompt; memory/retrieval; tool/schema/routing;
      output middleware; control flow/skill routing \\
    Causal stage & selection/activation; adherence/conflict;
      behavior/transformation; output/verification contract \\
    Structure & sparse or dense prevalence; single or interacting/masked
      double fault; correlated non-causal distractor present or absent \\
    Distribution & IID gate or held-out template/source shift \\
    Repairability & repairable; oracle-optimal no-fault; out-of-policy \\
    \bottomrule
  \end{tabular}
\end{table}

The generator freezes each oracle-correct harness, injected fault, true causal
surface and trigger, admissible repair, exact revert, neutral placebo,
protected-clean behavior, repairability, split groups, and artifact hashes.
The sampling target---not the observed outcomes---sets the factor weights.
Each independent family receives one preregistered primary optimizer seed,
shared across policies.  Family-level false promotion equals one if any invalid
edit is committed during that complete family--seed campaign.  Policies are
paired within family, and family-clustered inference integrates optimizer
randomness over the prospective seed draw.  A harmful promotion is a
commit whose atomic sealed audit fails the identical utility-benefit,
protected-behavior, or cost rule.  Accepted-edit precision counts only commits
that satisfy the oracle mechanism and every sealed gate; if its denominator is
zero it is reported as undefined, not 100\%.  Search seeds, scenario roles, and
rollouts within one generated family do not increase the family sample size.
Any secondary multi-seed analysis freezes either a family mean or an any-over-
seeds campaign event before outcomes and never treats seeds as independent
families.

\section{Factorial Decomposition and Fresh-Block Validity}
\label{app:factorial}

\begin{figure}[h]
  \centering
  \resizebox{0.88\linewidth}{!}{%
\begin{tikzpicture}[
  font=\sffamily,
  >=Latex,
  cell/.style={draw=black, rounded corners=2pt, align=center,
               minimum width=36mm, minimum height=14mm, inner sep=4pt},
  contrast/.style={-Latex, line width=0.8pt},
  fullcontrast/.style={-Latex, dashed, line width=0.9pt},
  heading/.style={font=\sffamily\bfseries, align=center},
  rowheading/.style={font=\sffamily\bfseries, align=right, text width=22mm},
  note/.style={font=\sffamily\footnotesize, align=center}
]
  \node[cell, fill=MutableFill] (ss) {\textbf{SS}\\score feedback\\performance rule};
  \node[cell, fill=TrustedFill, right=22mm of ss] (ms)
        {\textbf{MS}\\mechanism feedback\\performance rule};
  \node[cell, fill=GateGreen, below=18mm of ss] (sm)
        {\textbf{SM}\\score feedback\\mechanism rule};
  \node[cell, fill=SealedFill, right=22mm of sm] (mm)
        {\textbf{MM}\\mechanism feedback\\mechanism rule};

  \node[heading, above=5mm of ss] {Score-only\\optimizer view};
  \node[heading, above=5mm of ms] {Mechanism-visible\\optimizer view};
  \node[rowheading, left=4mm of ss] {Performance\\promotion};
  \node[rowheading, left=4mm of sm] {Mechanism\\promotion};

  \draw[contrast] (ss.east) -- node[above, note] {guidance} (ms.west);
  \draw[contrast] (ss.south) -- node[right, note] {verification} (sm.north);
  \draw[contrast] (sm.east) -- (mm.west);
  \draw[contrast] (ms.south) -- (mm.north);
  \draw[fullcontrast] (ss.south east) --
        node[sloped, below, note] {end-to-end} (mm.north west);

  \node[note, below=6mm of sm, text width=98mm, xshift=29mm]
        {Interaction: $(\mathrm{MM}-\mathrm{SM})-(\mathrm{MS}-\mathrm{SS})$.
         All cells retain the same ex-ante coordinate rules, authority calls,
         ceilings, and fresh-block schedule; treatment-induced search paths may diverge.};
\end{tikzpicture}%
}

  \caption{Controlled-fault factorial decomposition.  Horizontal contrasts
  isolate mechanism evidence as search guidance; vertical contrasts isolate
  mechanism-gated promotion; the dashed diagonal is the deployable end-to-end
  contrast.  Matching fixes ex-ante rules and budgets, not treatment-induced
  candidate contents.}
  \label{fig:factorial-decomposition}
\end{figure}

Let \(\mathcal F_{q-1}\) contain the joint, feedback-isolated history before
nomination batch \(q\).  Every candidate in that batch and every threshold
\(\alpha_{qaj}\) for condition \(a\) and gate dimension \(j\) is
\(\mathcal F_{q-1}\)-measurable before the fresh block is opened.  Assume that
for each true null,
\[
  \Pr(P_{qaj}\leq x\mid\mathcal F_{q-1})\leq x
  \quad\text{for all }x\in[0,1],
  \qquad
  \sum_{q,a,j}\alpha_{qaj}\leq\alpha.
\]
Thus a shorthand budget such as \(\alpha/Q\) can only be the total envelope
for nomination block \(q\); it must still be divided or validly combined across
conditions and gate dimensions so that the displayed global sum is at most
0.05.
Define a \emph{statistical false authorization} as a promotion for which at
least one prespecified fresh-block gate null is true.  Under the displayed
conditions, the campaign-wise probability of at least one such authorization
is at most \(\alpha\).  Indeed, that intersection-gate authorization rejects at
least one true null.  A union bound over those rejections, followed by
conditional super-uniformity and iterated expectation, gives
\(
  \Pr(\text{any statistical false authorization})
  \leq\sum_{q,a,j}\mathbb E[
  \Pr(P_{qaj}\leq\alpha_{qaj}\mid\mathcal F_{q-1})]
  \leq\sum_{q,a,j}\alpha_{qaj}\leq\alpha.
\)
Dependence from paired conditions is allowed.  Unique block identifiers alone
do not establish the premise: unopened source/family groups, candidate freezing,
isolated feedback, grader integrity, and complete consumption accounting are
required evidence.  Reusing a gate family after adapting to its aggregate
invalidates this argument even if the rollout seed changes.  This guarantee is
conditional on fresh blocks, calibrated tests, and complete alpha accounting.
It does not cover harm revealed only under sealed distribution shift, grader or
policy exploits outside the tested claims, provenance violations, or judge and
model misspecification; those are reported separately as harmful or protocol-
invalid promotions.  In particular, it is not a proof that the broader false-
promotion event in Section~\ref{sec:problem} is bounded by \(\alpha\).

\section{Figure Accessibility}
\label{app:figure-accessibility}

Figure~\ref{fig:mechanism-boundary} encodes authority by labeled containers and
line style, not color alone.  Left-to-right boxes in the upper container show
the mutable observe--diagnose--mutate loop.  The lower blue container shows a
matched quartet followed by activation, adherence, behavior, and
utility/protection/cost gates.  A dashed return arrow carries only a redacted
decision.  A separate tan box receives the commit registry and champion hashes
after all search cells close and releases one aggregate evaluation.  Candidate execution sits in a
dashed untrusted-runtime box without grader, target, or secret access.

Figure~\ref{fig:factorial-decomposition} is a labeled two-by-two grid.  Columns
change the optimizer view from score-only to mechanism-visible; rows change the
promotion rule from performance-only to mechanism-gated.  Solid horizontal and
vertical arrows denote guidance and verification contrasts.  A labeled dashed
diagonal from SS to MM denotes the end-to-end protocol contrast.  Labels, line
style, and position preserve all meaning without color.
